\section{先分清处理器能力、运行模式和指令宽度}

\topic{“64 位处理器”不表示上电第一条就是 64 位指令}
先直接回答最容易混淆的三个问题。

\begin{enumerate}[leftmargin=2.2em]
  \item \textbf{32 位或 64 位处理器怎样进入实模式？}
    它不是先以 32 位或 64 位方式运行，再执行一条指令进入实模式。硬件复位
    直接把控制寄存器和代码段状态设成实地址模式，软件看到的
    \code{CR0.PE=0}、\code{CR0.PG=0}。第一条软件指令已经处在实模式中。
  \item \textbf{复位后的指令是 16 位还是 32 位？}
    当前代码段的默认操作数宽度和默认地址宽度都是 16 位，所以没有宽度前缀
    的普通指令先按 16 位默认方式解码。
  \item \textbf{实模式里能不能使用 \code{EAX}？}
    80386 及后续处理器可以。在 16 位默认代码中，指令前缀可以把某一条指令
    的操作数改成 32 位。那一条指令会读写 \code{EAX}，处理器模式并没有
    随之改变。
\end{enumerate}

“位数”至少可能在说下面五件事。把它们压成一句“现在是 16 位”会丢掉必要
信息。

\begin{osgridlongtable}{L{0.19\linewidth}L{0.25\linewidth}L{0.27\linewidth}L{0.19\linewidth}}
名称 & 它回答的问题 & 本项目刚复位时的答案 & 决定者 \\ \hline
\endfirsthead
名称 & 它回答的问题 & 本项目刚复位时的答案 & 决定者 \\ \hline
\endhead
处理器能力 & 这颗 CPU 实现了哪些寄存器、指令和模式？ &
支持 x86-64，也兼容 16/32 位执行 & CPU 型号与 ISA \\ \hline
当前运行模式 & CPU 此刻按哪套分段、保护和分页规则工作？ &
实地址模式，PE=0、PG=0 & 控制寄存器、EFER 与段状态 \\ \hline
代码段默认宽度 & 没有宽度前缀时，操作数和地址按多少位解码？ &
默认操作数 16 位、默认地址 16 位 & 当前 CS 的隐藏属性 \\ \hline
本条指令宽度 & 这一条指令实际读写 AX 还是 EAX？ &
由操作码、默认值和前缀共同决定 & 指令字节 \\ \hline
汇编器编码策略 & NASM 应按哪种默认值选择机器码？ &
\code{bits 16} & 源文件中的汇编指示 \\ \hline
\end{osgridlongtable}

\code{bits 16} 不会被写成一条供 CPU 执行的指令。它只告诉 NASM：“假定
CPU 将按 16 位默认值解释后续字节；如果源码要求另一种宽度，请生成前缀。”
项目 Stage 1 的下面三条指令都在实模式中：

\begin{osgridlongtable}{L{0.24\linewidth}L{0.24\linewidth}L{0.20\linewidth}L{0.22\linewidth}}
源码 & 最终机器字节 & 实际操作宽度 & 执行后的模式 \\ \hline
\endfirsthead
源码 & 最终机器字节 & 实际操作宽度 & 执行后的模式 \\ \hline
\endhead
\code{mov eax, cr0} & \code{0F 20 C0} & 控制寄存器传送规定的宽度 & 实模式 \\ \hline
\code{or eax, 1} & \code{66 83 C8 01} & 32 位；\code{66} 覆盖 16 位默认值 & 实模式 \\ \hline
\code{mov cr0, eax} & \code{0F 22 C0} & 写 CR0；这一条把 PE 置为 1 & 进入保护模式转换点 \\ \hline
\end{osgridlongtable}

第二行是最好的反例：使用 \code{EAX} 不等于进入了“32 位模式”。执行
\code{or eax, 1} 时，改变的只有 EAX 和算术标志。直到第三行把候选值写入
CR0，\code{PE} 才真的改变。

\begin{pdfdiagram}
  \node[os hardware] (reset) {硬件 RESET};
  \node[os state,right=of reset] (real) {实模式\\默认 16/16};
  \node[os block,right=of real] (wide) {带 \code{66} 的一条指令\\32 位操作数};
  \node[os state,below=of wide] (realagain) {仍是实模式\\PE=0};
  \node[os block,left=of realagain] (pe) {\code{mov cr0, eax}\\PE 置 1};
  \node[os state,left=of pe] (protected) {far jump 后\\32 位代码段};
  \draw[os arrow] (reset) -- (real);
  \draw[os arrow] (real) -- (wide);
  \draw[os arrow] (wide) -- (realagain);
  \draw[os arrow] (realagain) -- (pe);
  \draw[os arrow] (pe) -- (protected);
\end{pdfdiagram}
\diagramnote{“32 位操作数”只描述一条指令读写的数据宽度；“32 位保护模式
代码段”描述后续指令的默认解码环境。两者不能互换。}

\section{复位状态、ROM 映射与首条指令}

\topic{复位是契约，不是初始化完成}
上电复位后，处理器处于实地址模式。最关键的可观察状态是指令入口：
\code{CS} 的可见值为 \code{0xF000}，隐藏基址为
\code{0xFFFF0000}，\code{IP/EIP} 指向 \code{0xFFF0}，因此第一条
指令的物理地址是：
\[
  0xFFFF0000 + 0xFFF0 = 0xFFFFFFF0.
\]
隐藏基址是理解现代 x86 复位地址的关键。此时不能简单套用普通实模式
\(segment \times 16 + offset\) 去解释第一次取指。

手册把寄存器记作 EIP，并不表示当前正在按 32 位默认方式执行。EIP 是处理器
已经实现的架构寄存器；当前代码段的默认宽度则是另一项状态。复位时
EIP=\code{0x0000FFF0}，本项目的 \code{E9 rel16} 只产生 16 位跳转目标，
所以我们在计算中写 IP=\code{0xFFF0} 和 IP=\code{0xF000}，是在明确只看
它此刻参与控制流的低 16 位。

\begin{pdfdiagram}
  \node[os hardware] (reset) {RESET};
  \node[os state,below left=of reset] (visible) {CS 可见值\\\code{0xF000}};
  \node[os block,below=of reset] (hidden) {CS 隐藏基址\\\code{0xFFFF0000}};
  \node[os block,below right=of reset] (offset) {EIP\\\code{0x0000FFF0}};
  \node[os hardware,below=of hidden] (physical) {首条物理地址\\\code{0xFFFFFFF0}};
  \node[os state,below=of physical] (reload) {实模式 far transfer 重载 CS\\
  base 变为 selector \(\times16\)};
  \draw[os arrow] (reset) -- (visible);
  \draw[os arrow] (reset) -- (hidden);
  \draw[os arrow] (reset) -- (offset);
  \draw[os arrow] (hidden) -- node[above left]{相加} (physical);
  \draw[os arrow] (offset) -- (physical);
  \draw[os arrow] (physical) -- (reload);
\end{pdfdiagram}
\diagramnote{依据 Intel SDM Vol. 3A 12.1.4 重绘。手册给出复位时 CS
可见值、隐藏基址和 EIP；项目再把结果映射到 128 KiB ROM 的文件偏移。}

\topic{128 KiB ROM 的顶端映射}
项目 ROM 覆盖 \code{0xFFFE0000} 到 \code{0xFFFFFFFF}。文件偏移与物理
地址之间满足：
\[
  physical = 0xFFFE0000 + file\_offset.
\]
所以复位向量的文件偏移是 \code{0x1FFF0}，入口
\code{0xFFFFF000} 的文件偏移是 \code{0x1F000}。

\begin{center}
\begin{tikzpicture}[os diagram]
  \node[os hardware,minimum width=42mm] (rom) {ROM 起点\\0xFFFE0000};
  \node[os state,above=12mm of rom,minimum width=42mm] (entry)
    {16 位入口\\0xFFFFF000};
  \node[os block,above=12mm of entry,minimum width=42mm] (reset)
    {复位向量\\0xFFFFFFF0};
  \node[os hardware,above=8mm of reset,minimum width=42mm] (end)
    {ROM 末端\\0xFFFFFFFF};
  \draw[os arrow] (rom) -- node[right]{128 KiB 地址空间} (entry);
  \draw[os arrow] (entry) -- (reset);
  \draw[os arrow] (reset) -- (end);
\end{tikzpicture}
\end{center}

\topic{三个复位字节 E9 0D F0 到底怎样跳转}
复位后的默认操作数宽度是 16 位，因此 CPU 读到 \code{E9} 后，还要读取
两个字节的 \code{rel16}。内存按低字节在前保存整数，所以
\code{0D F0} 合成 \code{0xF00D}。把它当作 16 位有符号数时：
\[
  \operatorname{signext}_{16}(\mathtt{0xF00D})=-4083.
\]

相对跳转不是从指令开头加位移。CPU 已经取完三个字节，顺序 IP 先从
\code{0xFFF0} 走到 \code{0xFFF3}，然后才加上 \(-4083\)：
\[
  0xF000 - (0xFFF0 + 3) = -4083,
  \qquad
  (-4083)\bmod 2^{16}=0xF00D.
\]
也可以正向验算：
\[
  \mathtt{0xFFF3}+(-4083)
  =\mathtt{0xF000}.
\]

这是 \emph{near jump}。它改写 IP，不重装 CS，也不写 CR0。因此跳转完成
以后，CS 可见值仍为 \code{0xF000}，CS 隐藏基址仍为
\code{0xFFFF0000}，处理器仍在实模式，默认宽度仍为 16 位。下一条指令的
物理地址才是：
\[
  \mathtt{0xFFFF0000}+\mathtt{0xF000}
  =\mathtt{0xFFFFF000}.
\]

\begin{osgridlongtable}{L{0.24\linewidth}L{0.24\linewidth}L{0.24\linewidth}L{0.18\linewidth}}
状态 & 跳转前 & 跳转后 & 是否改变 \\ \hline
\endfirsthead
状态 & 跳转前 & 跳转后 & 是否改变 \\ \hline
\endhead
IP/EIP 的有效低 16 位 & \code{0xFFF0} & \code{0xF000} & 是 \\ \hline
CS 可见值 & \code{0xF000} & \code{0xF000} & 否 \\ \hline
CS 隐藏基址 & \code{0xFFFF0000} & \code{0xFFFF0000} & 否 \\ \hline
CR0.PE / CR0.PG & 0 / 0 & 0 / 0 & 否 \\ \hline
当前模式 & 实模式 & 实模式 & 否 \\ \hline
默认操作数/地址宽度 & 16 / 16 & 16 / 16 & 否 \\ \hline
\end{osgridlongtable}

\topic{near、far 和模式切换不是同一个动作}
\code{near jump} 只在当前代码段内改变指令偏移。\code{far jump} 同时
给出新的 CS 和偏移，所以会重装 CS 的隐藏部分。但 far jump 本身也不等于
“进入保护模式”：实模式中的 far jump 仍然跳到实模式代码，只有
\code{CR0.PE} 改变以后，处理器才按保护模式规则用选择子查 GDT。

\begin{osgridlongtable}{L{0.20\linewidth}L{0.23\linewidth}L{0.24\linewidth}L{0.23\linewidth}}
动作 & 改写的主要状态 & 会不会单独改模式 & 本项目中的用途 \\ \hline
\endfirsthead
动作 & 改写的主要状态 & 会不会单独改模式 & 本项目中的用途 \\ \hline
\endhead
\code{E9 rel16} near jump & IP & 不会 & 从复位向量跳到同一特殊 CS 下的 ROM 入口 \\ \hline
实模式 far jump / \code{retf} & CS:IP 及 CS 隐藏基址 & 不会 & ROM 把控制权交给 RAM 中的 Stage 1 \\ \hline
写 \code{CR0.PE=1} & 保护模式使能位 & 会，进入转换点 & 启用保护模式规则 \\ \hline
PE 置位后的 far jump & CS 与代码段属性 & 模式已由 PE 改变 & 装入 GDT 的 32 位代码段并建立新默认宽度 \\ \hline
\end{osgridlongtable}

\begin{failurebox}
如果使用会重载 CS 的远跳转，就必须重新证明目标段如何映射到 ROM；如果链接
脚本把入口放错一个字节，CPU 也不会给出友好错误，只会执行填充值或触发异常。
\end{failurebox}

\topic{复位向量固定在末端的历史连续性}
8086 只有 20 位物理地址总线，复位后从物理地址 \code{0xFFFF0} 开始执行，
这个位置距 1 MiB 地址空间末端只有 16 字节。把入口放在顶端，使主板能够把
固件映射在高地址，同时把大部分低地址留给 RAM。后续处理器扩展了地址宽度，
但软件生态已经依赖这一入口语义。现代 x86 在复位时借助 CS 的可见选择子和
隐藏基址组合，把第一条取指定位到 \code{0xFFFFFFF0}，继续保持“距 4 GiB
末端 16 字节”的结构。

兼容性使新处理器可以运行既有固件，也使现代启动过程必须从一个极受约束的
历史状态开始。AMD64 没有把复位入口直接改成长模式，因为这会同时破坏固件
生态、主板设计和诊断工具。长模式因此是软件逐步建立的目标状态，而不是上电
默认状态。

\topic{近跳转字节的逐项推导}
复位位置只有 16 字节可供使用，真正的初始化代码必须放在 ROM 的其他位置。
\code{JMP rel16} 的机器编码由一个字节操作码和一个小端序的 16 位有符号
位移组成。处理器先把 IP 推进到下一条指令，再把位移加到 IP：
\[
  \mathrm{IP}_{next}=0xFFF0+3=0xFFF3,
  \qquad
  \mathrm{IP}_{target}=0xFFF3+\operatorname{signext}(0xF00D)=0xF000.
\]

\begin{osgridlongtable}{L{0.16\textwidth}L{0.22\textwidth}L{0.48\textwidth}}
\textbf{字节位置} & \textbf{数值} & \textbf{处理器解释}\\ \hline
\endfirsthead
\textbf{字节位置} & \textbf{数值} & \textbf{处理器解释}\\ \hline
\endhead
\code{0xFFFFFFF0} & \code{E9} & 选择近相对跳转指令\\ \hline
\code{0xFFFFFFF1} & \code{0D} & 位移低字节\\ \hline
\code{0xFFFFFFF2} & \code{F0} & 位移高字节；组合为 \code{0xF00D}\\ \hline
\code{0xFFFFFFF3} & \code{FF} & 不属于该指令；正常执行不会到达\\ \hline
\end{osgridlongtable}

近跳转不写 CS，因此隐藏基址保持 \code{0xFFFF0000}。新的线性地址为
\code{0xFFFF0000 + 0xF000 = 0xFFFFF000}，正好落在链接器安排的 ROM
入口。这个推导同时连接了机器码、小端序、符号扩展、指令长度、段缓存和链接
布局，任何一个量变化都必须重新证明。

\topic{填充值也是系统行为}
ROM 未占用区域填充为 \code{0xFF}，能够使镜像大小和烧录语义保持确定，但
\code{0xFF} 并不是安全的“空指令”。控制流意外进入填充区后，处理器可能将
连续字节解释为有效指令前缀与操作码，也可能触发异常。早期尚未建立异常处理，
最终表现通常只是复位、停机或无输出。因此链接断言和二进制字节检查必须在
运行前阻止布局错误。
